\section{Fixed-target comparison of contact graphs}
\label{sec:embedded-comparison-v163}
This section makes the change of category in the contact construction
explicit. A formal congruence depending on the pencil parameter acts
on an incidence bundle; it is not an automorphism of the fixed target
whose Hilbert scheme is under consideration. We retain the incidence
projection and undo every such change before taking a Hilbert point.

\subsection{Objects, targets, and the gluing statement}
Work on a finite-type framed atlas $S$ of regular pencils avoiding
corank at least three, with central pencil $R_0$. A frame here means
a lift of the projective source, an ordered pencil basis, and a local
parameter at each of the finitely many central corank-two points.
These choices are not quotiented in an embedded Hilbert fibre. Let
$\mathcal L\subset\Pj(\Sym^2V)\times S$ be the universal pencil
line, $X=\mathrm{CQ}(V)$, and
\[
 \mathcal Y=\mathcal L\times_{\Pj(\Sym^2V)}X.
\]
The targets and their maps are
\begin{equation}\label{eq:incidence-diagram-v163}
 \begin{array}{ccccc}
 \mathcal Z&\lhook\joinrel\longrightarrow&\mathcal Y
             &\lhook\joinrel\longrightarrow&X\times S\\
 &&\downarrow&&\downarrow\\
 &&\mathcal L&\longrightarrow&S .
 \end{array}
\end{equation}
Here $\mathcal Z$ denotes an embedded flat family, and the last
vertical arrow is projection. The middle map is a closed immersion
because $\mathcal L$ is the universal family of actual lines.
An object in the graph retains its point of $S$, not just the curve
in $X$. On the complement of the central contact supports, the
complete-quadric lift is the forced section of this incidence target.

Write $\widehat S=\operatorname{Spf}\widehat\OO_{S,R_0}$.
For a local Artin complex algebra $T$ with a specified map to
$\widehat S$, define $\mathcal D_X(T)$ to be the set of embedded
$T$-flat closed subschemes of \eqref{eq:incidence-diagram-v163},
with the fixed central curve, the specified Hilbert polynomial, and
the forced lift off the contacts. Isomorphism means equality as
closed subschemes, with the identity on the retained base. Restrict
this functor, when needed, to the completed reduced closure of the
nonincident graph. We do not use the abstract deformation functor
of the curve or identify objects by an ambient automorphism.

\begin{lemma}[Patching closed ideals with their quotients]
\label{lem:ideal-patching-v163}
Let $T$ be a local Artin complex algebra, let $W$ be a Noetherian
ambient scheme over a neighbourhood of a point of $\mathcal L_T$,
and complete along that point's inverse image. Let $W^\circ$ be
its complement. Compatible coherent ideal quotients on $W^\circ$
and the completion, with their identification on the punctured
completion, glue uniquely to a coherent ideal quotient on $W$.
The assertion applies to quotients with torsion supported at the
contact and to their flat families over $T$. Flatness, equality of
ideals and the specified closed incidence equations can be checked
on this cover. The same statement holds for finitely many disjoint
contact supports and after any Artin base change.
\end{lemma}
\begin{proof}
Use the ring-theoretic completion on affine neighbourhoods and its
pullback on projective charts. The map from this completion is flat,
Noetherian, and is an isomorphism over the closed support. The
formal patching equivalence for quasi-coherent sheaves is therefore
\cite[Proposition~5.6(4), Example~5.8]{Bhatt163}; its hypotheses are
also covered by Theorem~1.4 and Lemma~5.12(2) there. Flatness of the
completion is essential: the unrestricted module assertion for a
nonflat completion is not the statement being invoked. In
particular no torsion-freeness at the contact is imposed.

Apply this symmetric monoidal equivalence to the quotient with its
unit and multiplication. The ambient structure map and the ideal
kernel are retained as part of the data. Surjectivity and the
ordinary degree-zero algebra property are detected on the jointly
surjective flat cover consisting of the completion and the open
complement. Thus the glued object is an ordinary ideal quotient.
Finite presentation descends on this faithfully flat cover; the
ambient rings are Noetherian. The same cover detects base flatness
by the local criterion for flatness, or by vanishing of the relevant
Tor modules. It detects all the additional incidence equations.
Patching on the projective affine cover then glues the coherent
quotients. Iterating over disjoint supports gives the last assertion.
This argument supplies descent data on the punctured completions;
faithfulness of completion alone would only give uniqueness, not
existence.
\end{proof}

At the $i$th contact split a nonsingular block of size $n-2$. The
remaining symmetric Schur complement has entries $a_i,b_i,c_i$.
On this rank open the only noninvertible minor ideal below the
determinant is
\[
                  J_i=(a_i,b_i,c_i)=I_{n-1}.
\]
The determinant factor is an invertible ideal on the integral total
space and does not change the simultaneous graph. A congruence
$M_i(x_i)$ and the Schur splitting replace this generating triple
by an invertible linear transformation over the completed base
ring. Consequently they give an isomorphism between the corresponding
projective graph targets over the formal parameter disc. Denote it
by $\Theta_i$; it preserves $x_i$ and is part of the comparison data.
In normal coordinates the triple is $(f_i,g_i,r_i)$, the fourth
polynomial $k_i$ having been removed by the invertible row operation
$c_i\mapsto c_i-k_i a_i$.

\begin{proposition}[Completed embedded comparison]
\label{prop:embedded-functors-v163}
Fix the framed atlas, formal Schur splittings and normalizing
congruences just described. Let $\mathcal D_{\rm nor}$ be the
functor of embedded normal-graph quotients on the formal contact
discs, together with their identification with the forced lift on
the punctured discs. Then there is an equivalence of functors
\begin{equation}\label{eq:functor-comparison-v163}
 \mathcal D_X\ \simeq\ \mathcal D_{\rm nor}
\end{equation}
obtained by restriction and $\Theta_i$, with inverse given by
$\Theta_i^{-1}$ and ideal patching. No congruence quotient is taken
on either side. Changing the chosen frames or congruences conjugates
these functors and their inverse maps; it does not identify distinct
fixed-target Hilbert points.

On the reduced graph component, the complete local rings are the
complete local rings of the corresponding polynomial graph charts,
with the smooth parameters of the contact atlas adjoined. In
particular this assertion holds both at the division charts of
Theorem~\ref{thm:division-chart-v162} and at every determinantal
conic chart of Lemma~\ref{lem:HB-family-v163}. The comparison is
compatible with restriction to fibres. Where these total charts
are regular, normalization of the total graph changes neither them
nor their scheme fibres.
\end{proposition}
\begin{proof}
The completed normal form of
Theorem~\ref{thm:ambient-versal-v162} is used over
$T[[x_i]]$, with the base maximal-ideal filtration for its inductive
construction. The parameter $x_i$ is fixed; its powers specify the
remainder degrees, not the topology in that induction. Since $T$
is Artin the base induction is finite. For the complete universal
base it is the inverse limit of these finite inductions, so no
interchange of two unrelated convergence assertions is required.

An embedded family restricts to a coherent quotient on each
completed contact neighbourhood and to the forced quotient on the
complement. Apply $\Theta_i$ on the first and the induced
identification on the punctured neighbourhood. Conversely undo
$\Theta_i$, impose the original incidence equations, and apply
Lemma~\ref{lem:ideal-patching-v163}. These two operations are inverse
on quotients and their structure maps, proving \eqref{eq:functor-comparison-v163}
for every $T$, including nonreduced bases. The use of the image Rees
algebra can be recorded explicitly: a homomorphism $R\to R'$ gives
\[
        \mathcal R_R(J)\otimes_RR'\twoheadrightarrow
                          \mathcal R_{R'}(JR').
\]
Proj of this surjection embeds the actual image graph into the
base-changed incidence target. We do not assert that this
surjection is always an isomorphism. The closed ideals compared
in the proposition are these actual embedded ideals, rather than
an arbitrary base change of a blow-up.

For precision about the graph component, work first on the reduced
finite-type Hilbert graph over the atlas. The closure ideal of its
nonincident open can be computed in each finite presentation by
saturating with the ideal of the excluded base locus. Noetherianity
makes the increasing sequence of ideal quotients stationary; flat
completion commutes with each finite ideal quotient. It therefore
commutes with this closure computation. The rings are excellent
complex algebras, and their reduced completions stay reduced.
The patching equivalence preserves the nonincident section and the
retained base, hence also this specified reduced closure. It is
not a statement that unrelated components of an abstract curve
Hilbert scheme have disappeared.

The contact map on the completed smooth framed base is formally
smooth, with the bounded coefficient parameters as independent
coordinates and the remaining directions as power-series variables.
The scalar $k_i$ is one of these independent smooth factors for
the graph problem: the row operation above removes it without
changing the embedded ideal when the operation is undone.
Different contact discs patch independently. The polynomial graph
charts give algebraic representatives for their contact quotients;
their recovery maps retain the base coefficients. Applying the
functor equivalence gives the asserted isomorphism of completed
local rings with the smooth factors adjoined.

This is not yet an algebraization by formal assertion. In each
application below the morphism on the special fibre is constructed
algebraically first: a primitive line over the finite contact
algebra, or a conic together with the fixed main curve, supplies
its explicit coherent ideal. The finitely many coefficients of
the central gauge on that contact algebra are algebraic. Thus
these are finite-presentation morphisms into the actual Hilbert
graph. The completed comparison proves their local isomorphism
on the specified fibre charts. The residue fields at closed
points are $\C$ and the local rings are Noetherian; the formal
criterion for etaleness, \cite[Lemma~41.11.3]{Stacks163}, applies.
The coefficient or ideal recovery gives a monomorphism. Such a
morphism is an open immersion on these neighbourhoods. This orders
the argument as algebraic construction, completed comparison, then
etale monomorphism, rather than treating formal normal coordinates
as an algebraic family without descent.

An alternative normalizing congruence acts through a transformation
of the same relative incidence target. Both directions of the
comparison compose with that transformation and its inverse. The
composition in the fixed target remains the identity. The same
identity is the cocycle on overlaps of the source-frame atlas, so
the comparison descends. Finally a regular total chart is normal;
the finite normalization of the total graph is its identity there.
Taking its fibre retains the original base equations, possibly
with nilpotents. No commutation of fibre and normalization is used.
\end{proof}

\subsection{Cohomological recovery along the primitive locus}
\begin{lemma}[The union sequence and the retained coefficient base]
\label{lem:recovery-details-v163}
Let $D=\Spec\C[x]/(x^a)$ and let $Z_w=C_0\cup T$, with
$T=\Pj^1_D$ the primitive tail and intersection the attaching
section $D$. For $L=\OO(a-1,1)|_{Z_w}$ there is an exact sequence
\begin{equation}\label{eq:union-cohomology-v163}
 0\to L\to L|_{C_0}\oplus L|_T\xrightarrow{\operatorname{res}_C-
 \operatorname{res}_T}L|_D\to0.
\end{equation}
The restriction from $T$ is onto, $H^1(Z_w,L)=0$, and
$h^0(Z_w,L)=2a$. For each such $w$ these assertions persist on a
neighbourhood in the integral graph $\Gamma_a$ over $B_a$. On that
neighbourhood the recovery map has a line-bundle cokernel, and
vanishing of the retained third coordinate in that cokernel is a
scheme-theoretic equality, not a fibrewise assertion.
\end{lemma}
\begin{proof}
The ideals of a scheme-theoretic union give
$0\to\OO_{Z_w}\to\OO_C\oplus\OO_T\to\OO_D\to0$; tensor by
$L$. A line bundle on the local Artin scheme $D$ is free, and the
evaluation $H^0(\Pj^1_D,\OO(1))=\OO_D^2\to\OO_D$ at its section
has a unit coefficient, hence is surjective. On the main component
$L=\OO(a-1)$ and on the tail it is $\OO(1)$ up to an invertible
Artin factor. The cohomology sequence gives $a+2a-a=2a$ and the
claimed vanishing, over the Artin ring rather than only its residue
field.

To specify homogenization, write $W$ for degree-$(a-1)$ homogenization
of $w$, and $Q_0$ for degree-$(a-2)$ homogenization of $q$; when
$a=1$ the latter space is zero. The map is
\[
 (W,Q_0)\longmapsto WG-sQ_0F
       \quad\hbox{in }H^0(Z,\OO(a-1,1)).
\]
The section to be recovered is $s^{a-1}R$. There are $a+(a-1)$
unknown coefficients. Restricting to the main component, then to
the tail, proves injectivity at every central $w$. Upper
semicontinuity and cohomology and base change give a pointwise open
neighbourhood on which the map is a split injection of ranks
$2a-1$ and $2a$. These neighbourhoods cover the whole central
locus; no single uniform affine shrink is asserted. Its cokernel
is a line bundle. The class of $s^{a-1}R$ is zero on the dense
nonincident graph. A section of a line bundle on an integral
reduced neighbourhood that vanishes densely is zero. The recovered
coefficients are thus regular and satisfy the division identity.
The projection to $B_a$ is indispensable: it already supplies
$f,g,r$. All comparisons retain it.
\end{proof}
